\documentclass[11pt]{report}

\usepackage{graphicx}
\usepackage{booktabs}
\usepackage{hyperref}
\usepackage{listings}
\usepackage{amsmath}
\usepackage{geometry}
\usepackage{tabularray}
\geometry{margin=1in}

\title{Micromouse Documentation}
\author{McMaster Micromouse}
\date{Last Updated: July 30, 2026}

\begin{document}

\maketitle
\tableofcontents
\clearpage

\chapter{Introduction}
    \section{What is Micromouse?}
	A Micromouse is a small autonomous robot that navigates and solves a maze as quickly as
	possible. The robot begins with no knowledge of the maze layout and must explore it using 
	external sensors. As it explores, it constructs an internal map of the maze, determines 
	the most efficient route, and then completes subsequent runs at increasingly high speeds.
	\newline

	\noindent
	Micromouse has existed as an international engineering competition hosted by IEEE for 
	decades, and remains a great project for learning embedded systems. A successful Micromouse
	requires the integration of mechanical design, electronics, and software development. 

    \section{Club Objectives}
	The goal of our club is to make the process of designing, building, and programming a 
	Micromouse accessible to students of all experience levels. Robotics projects often require 
	knowledge from multiple disciplines, which can make them difficult 	to approach independently. 
	Through structured workshops, hands-on resources, and comprehensive documentation, the club 
	aims to lower this barrier and provide a clear path from beginner to autonomous robot.
	\newline

	\noindent
	Participants will be guided through the complete development process, from assembling hardware
	and programming embedded systems to implementing sensing, localization, and path-planning 
	algorithms. Rather than simply providing a finished design, the club emphasizes understanding 
	the engineering principles behind each subsystem so members can modify, improve, 
	and extend their own robots.
	\newline

	\noindent
	While the club provides a complete reference robot, members are encouraged to explore their own 
	ideas and improvements. Topics such as custom PCB design, alternative hardware, and advanced 
	control strategies are beyond the scope of the workshops but are left open for exploration. 
	The provided platform should be viewed as a foundation upon which participants can experiment, 
	iterate, and develop their own designs.

    \section{How to Use This Guide}
	This document serves as both a workshop and technical reference for the Micromouse Club. It is 
	intended to complement our workshops by providing detailed explanation of concepts, design decisions, 
	and implementation details that we may not have time to cover during workshops. It also allows 
	members to revisit or look ahead. 
	\newline

	\noindent
	The guide is organized into four main sections:
	\begin{enumerate}
		\item \textbf{Hardware:} Introduces the electrical components used in the robot. Explains how
				they operate and covers wiring.
		\item \textbf{Software:} Describes sensor integration, motion control, localization, and path-
				finding algorithms. 
		\item \textbf{Reference Implementation:} Documents the club's reference Micromouse design. This
				section serves as a starting point for members building their own robots.
	\end{enumerate}

\chapter{System Overview} \label{ch:system_overview}
    \section{Robot Architecture}
    \section{Hardware Overview}
    \section{Software Overview}

\part{Hardware}

\chapter{Materials} \label{ch:materials}
    \section{Components} \label{sec:components}

	There are a variety of components that can be used for a Micromouse. The main categories are summarized
	in \nameref{ch:system_overview}. Table \ref{tab:component_list} lists all the components provided in 
	our kits.
	\newline

	\begin{table}[ht]
		\centering
		\begin{tblr}{
			vline{1-4} = {-}{},
			hline{-} = {1-3}{},
			} 
			\textbf{Components} & \textbf{Quantity} & \textbf{Purpose} 						\\ 
			ESP32-DEVKIT		& 1 & Selected microcontroller. Main "brain" of the bot.	\\ 
			VL53L1X    			& 3 & Time of Flight sensor. Used to align the bot.        	\\ 
			TB6612FNG  			& 1 & Motor driver; to control the motors.       			\\ 
			G12-N20 Motors  	& 2 & Motors for the wheels.								\\ 
			N20 Motor Encoders  & 2 & Attached to motors to track of motor position.		\\ 
			N20 Gear Wheels  	& 2 & Wheels for the bot to move.							\\ 

		\end{tblr}
		\caption{Component List}
		\label{tab:component_list}
	\end{table}
	
	\noindent
	These provide a good starting point for a basic Micromouse bot, and will be the parts we focus on 
	throughout our workshops. The \nameref{ch:hardware_theory} section will also go in-depth on how each
	component works. However, you are free to choose your own parts. Some alternative options are 
	described in the next section.

    \section{Alternatives} \label{sec:alternatives}
	This section briefly introduces several alternative components that can be used for your own 
	Micromouse, along with the rationale behind our selected components. Keep in mind
	that the options presented are not exhaustive; you are encouraged to research additional 
	alternatives. A summary list (Table \ref{tab:component_alternatives}) is included at the end of this section.
	\newline

	\noindent
	\textbf{Microcontroller}
	\newline
	Micromouse robots require enough of computing power, memory, and peripheral support to handle 
	sensor readings, motor control, and path-finding algorithms. While basic Arduino boards such as the  
	Arduino Uno or Nano(ATmega328P) can be used for simple robots, they become restrictive for more advanced
	algorithms.
	\newline

	\noindent
	The STM32 family (e.g. the STM32F103, 'blue pill, 'black pill', and STM32F4 Series) is a popular choice for 
	more advanced Micromouse bots due to its performance, low-level hardware control, and extensive timer and 
	interrupt capablities. However, due to the learning curve of using the STM32CubeIDE, can be more challenging
	for beginners. For this reason, an ESP32 board was chosen to provide a more accessible development while 
	still providing sufficient performance for a Micromouse.
	\newline

	\noindent
	Alternative ESP32 boards, such as the ESP32-C3 and ESP32-S3, can also be considered. These boards 
	are smaller and have different peripheral options, but have fewer available GPIO pins compared 
	to some ESP32 development boards. With careful planning, they can still support the required 
	sensors, motor drivers, and encoders.
	\newline

	\noindent
	\textbf{Distance Sensors}
	\newline
	Distance sensors are one of the most important parts of a Micromouse bot. When selecting a sensor,
	it is important to consider range, update rate, field of view, and physical size. 
	\newline
	
	\noindent
	We will be using the VL53L1X time-of-flight (ToF) sensor due to its relatively long sensing range,
	accuracy, and size. Alternatives of the ToF sensor include the VL53L0X (shorter range) or VL53L4CX
	(longer range). 
	\newline

	\noindent
	Infrared (IR) distance sensors are another common choice, especially in advanced Micromouses. 
	Compared to ToF sensors, IR sensors have faster update rates and lower latency. However, they 
	require	analog signal processing and require more calibration.  
	\newline

	\noindent
	Ultrasonic sensors may also be used for slower bots, but are not recommended in general. They 
	are physically larger and are slower since they use a sound waves rather than light.
	\newline

	\noindent
	\textbf{IMU Sensors}
	\newline
	For advanced and faster Micromouses, it is worth considering adding an IMU (inertial measuring
	unit). Common choices include the MPU6050 and MPU6500. These provide more accurate turns and 
	smoother motion control, but are not strictly necessary for beginner robots. As a result, we 
	have not provided any in the kits as beginner robots work with just wheel encoders.
	\newline

	\noindent
	\textbf{Motors}
	\newline
	Micromouse motors do not need huge amounts of power. The focus is on speed, control, smoothness,
	size, and accurate encoder readings. Encoders often come together with the motor.
	\newline

	\noindent
	N20 DC Gear Motor would work are common choices for beginner Micromous bots. They come in various 
	gear ratios, are compact, cheap and widely available. The ones we provide have a gear ratio of 
	30:1, however other ratios can also be used. Coreless DC gear motors are also worth looking into 
	for advanced Micromouses.
	\newline

	\noindent
	Stepper motors are not recommended for Micromouse bots. Although they provide precise position
	control, they are larger, heavier, consume more power, and have lower speeds and
	acceleration compared to DC motors.
	\newline

	\noindent 
	\textbf{Motor Driver}
	\newline
	The motor driver acts as the bridge between the microcontroller and the motors. Since we are 
	just using N20 gear motors, the chosen motor driver was the TB6612FNG. It is a popular choice 
	for small bots as it is small, cheap and can control two separate motors. 
	\newline

	\noindent
	Alternatives include the DRV8833 (similar to TB6612FNG) or the DRV8871 (for larger motors). 
	The L298N is not really recommended due to its bulky size and large voltage drops. 
	\newline

	\noindent
	\textbf{Summary}
	\begin{table}[ht]
		\centering
		\begin{tblr}{
			vline{1-5} = {-}{},
			hline{-} = {1-4}{},
			} 
			\textbf{Components} & \textbf{Selected} & \textbf{Alternative} 					 \\ 
			Microcontroller		& ESP32-WROOM 			  & ESP32-C3, ESP32-S3, STM32F103/F4 \\ 
			Distance Sensor    	& VL53L1X 				  & VL53L0X, VL53L4CX, IR sensors    \\ 
			IMU   				& None 					  & MPU6050, MPU6500      			 \\ 
			Motors  			& N20 DC Gear Motor(30:1) & Other N20 gear ratios, Coreless DC motors \\
			Motor Driver  		& TB6612FNG				  & DRV8833, DRV8871						\\ 

		\end{tblr}
		\caption{Summary of Component Alternatives}
		\label{tab:component_alternatives}
	\end{table}

\chapter{Hardware Theory} \label{ch:hardware_theory}
    \section{ESP32}
    \section{Motors}
    \section{Motor Driver}
    \section{Encoders}
    \section{Time-of-Flight Sensors}
    \section{Power System}

\chapter{Electrical Design}
    \section{Wiring Schematic}
    \section{Pin Assignments}

\part{Software}

\chapter{Software Architecture}
    \section{Project Structure}
    \section{Control Loop}

\chapter{Firmware Setup}
    \section{Development Environment}
    \section{Libraries}
    \section{Uploading Code}

\chapter{Sensors}
    \section{Reading ToF Sensors}
    \section{Reading Encoders}
    \section{Sensor Calibration}

\chapter{Motion Control}
    \section{Differential Drive}
    \section{PID Control}

\chapter{Mapping and Localization}
    \section{Maze Representation}
    \section{Localization}
    \section{Wall Detection}

\chapter{Path Planning}
    \section{Flood Fill}
    \section{Shortest Path Optimization}
    \section{Alternative Algorithms}

\part{Reference Implementation}

\chapter{Mechanical Assembly}
\chapter{Electronics Assembly}
\chapter{Firmware Configuration}
\chapter{Testing}

\appendix

\chapter{Pinout Tables}
\chapter{Datasheets}
\chapter{Equations}
\chapter{Troubleshooting}
\chapter{Glossary}

\end{document}